\documentclass[10pt,a4paper]{amsart}
\usepackage[margin=19mm]{geometry}
\usepackage{amsmath,amssymb,mathtools}
\usepackage[scr=rsfs,cal=euler]{mathalfa}
\usepackage{xfrac}
\usepackage[unicode,hidelinks,bookmarks=false]{hyperref}
\newcommand{\ie}{i.e.\ }
\newcommand{\cf}{cf.\ }
\newcommand{\resp}{respectively\ }
\newcommand{\Subsection}{\subsection}
\providecommand{\nobreakdash}{\nobreak\hbox{-}}
\setlength{\emergencystretch}{2em}
\multlinegap0pt
\usepackage{amssymb}
\usepackage{xcolor}
\newtheorem{lemma}{Lemma}
\title{The multivariate Hermite method for counting real and complex solutions to polynomial systems}
\author{Volodymyr Oleksiyuk}
\date{Source: arXiv:2510.23897v2 (CC BY 4.0)}
\begin{document}
\maketitle

\noindent\emph{Source and licence.} The multivariate Hermite method for counting real and complex solutions to polynomial systems. Version arXiv:2510.23897v2 (CC BY 4.0), distributed by arXiv under the Creative Commons Attribution 4.0 International licence, http://creativecommons.org/licenses/by/4.0/, as recorded in the arXiv licence metadata for this exact version and shown on the abstract page https://arxiv.org/abs/2510.23897v2. Full licence notice: \texttt{licenses/arxiv-2510.23897-CC-BY-4.0.txt}.\par\smallskip

\noindent\textit{Editorial scope.} Counted unit: the article's lemma lemma (source0.tex lines 298--309). The article's complete original proof of it is delivered with it (source0.tex lines 311--334, one \texttt{proof} environment of 24 lines). No mathematics is written, repaired or re-derived: the statement and the proof are the article's own lines, in the article's own notation, and no retained line is substituted or re-flowed.  No further source-local context is needed: every cross-reference of every retained line resolves inside the two delivered spans.  Nothing was omitted from any retained span, so the package carries no omission record.  No retained line contains a citation, so the wrapper carries no bibliography and no bibliography entry.  No retained line contains a figure environment, an image-inclusion command or drawing code, so no image is delivered and the package declares no asset dependency.  Editorial formatting only, in addition: an AMS \texttt{amsart} preamble that restates the article's own declarations and macros used by the retained lines, namely the environments \texttt{lemma} on separate counters, together with the standard packages those lines need.  The retained environments print in this document as Lemma 1 in order of appearance, whereas the article prints the same items with the numbering of its own full document.  The complete native source of the article as distributed in the arXiv e-print for this exact version is bundled unchanged as \texttt{sources/187913/source0.tex}.
\begin{lemma}
For any ideal $I \triangleleft \mathbb{R}[X_1, \ldots , X_n]$ and $A = \mathbb{R}[X_1, \ldots , X_n]/I$ we have
\[
\dim(A/\sqrt{0})<\infty \Leftrightarrow \dim(A)<\infty
\]
as vector spaces, where 
\begin{align*}
    \sqrt{0}&:=\{a \in A \, | \, \exists n \in \mathbb{N}: a^n = 0\} \\
    &= \{ p+I \, | \, p \in \mathbb{R}[X_1, \ldots , X_n], \exists n \in \mathbb{N}: \, p^n \in I\} = \sqrt{I}/I
\end{align*}
is the nilradical over $A$.
\end{lemma}

\begin{proof}
As $\sqrt{0}$ is a linear subspace of $A$, let us consider the linear decomposition of $A$. Let $U$ be the linear complement of $\sqrt{0}$, i.e.~$A = \sqrt{0} \oplus U$. 

Note that $U \cong A/\sqrt{0}$ as linear subspaces, under the residue isomorphism. It follows that $\dim(A) < \infty$ if and only if both $\dim(A/\sqrt{0}) < \infty$ and $\dim(\sqrt{0}) < \infty$. Thus, it suffices to prove that $\dim(\sqrt{0}) < \infty$ under the assumption $\dim(A/\sqrt{0}) < \infty$.

Since $A$ is a finitely generated algebra over a field, it is Noetherian. Therefore, the ideal $\sqrt{0} \triangleleft A$ is finitely generated as an ideal.

Let $\sqrt{0} = \langle g_1, \ldots, g_k\rangle$. Let $N_j \in \mathbb{N}$ be the smallest natural numbers such that $g_j^{N_j} = 0$ for all $j \in \{1, \ldots, k\}$, and let $N = \max\{N_1, \ldots, N_k\}$. Then for all $h \in \sqrt{0}$ we have $h = \sum_{j=1}^k h_j g_j$ with $h_j \in A$, and
\[
h^M = \sum_{|\alpha|=M} \widehat{h_\alpha} g^\alpha, \quad \forall M \in \mathbb{N},
\]
where $\alpha$ is a non-negative multiindex. This means however, that for $M \geq k(N-1)+1$ we have $h^M = 0$ by the pigeonhole principle.

Now for every $h \in \sqrt{0}$ with $h = \sum_{j=1}^k h_j g_j$, $h_j \in A$, we can decompose
\[
\forall j \in \{1, \ldots, k\}: \quad h_j \in A \;\Rightarrow\; h_j = \widetilde{h_j} + h_j', \quad
\widetilde{h_j} \in U, \quad h_j' \in \sqrt{0}.
\]
We can now take each of the $h_j'$ as a new $h$ and again decompose it over our ideal. Continuing this process, after at most $M$ decompositions we reach products of $g_j$ of total degree $M$, i.e.~a point where our original $h$ is decomposed as 
\[
h = \sum_{|\alpha| < M} \widehat{h_\alpha} g^{\alpha}, \quad \widehat{h_\alpha} \in U,
\]
since all terms in $\sqrt{0}$ will satisfy $|\alpha| \geq M$. This sum has finitely many terms. We are working under the assumption that $A/\sqrt{0}$, and hence $U$, is finite-dimensional. Therefore, decomposing the $\widehat{h_\alpha}$ in $U$ yields a finite linear combination, and thus $\dim(\sqrt{0}) < \infty$.
\end{proof}

\end{document}
